% ============================================================================
%  ESQUELETO DE PRESENTACIÓN
%  Informe: report/main.pdf
%
%  Tema: Boadilla + color theme "whale" (Beamer)
%
%  Compilar desde esta carpeta:
%      latexmk -pdf main.tex
%  o bien:
%      pdflatex main.tex && pdflatex main.tex
%
%  Las figuras del informe están en ../report/Figures/.
%  Cada "% TODO" indica un hueco que hay que completar.
% ============================================================================

\documentclass[aspectratio=169,11pt]{beamer}

% ---------------------------------------------------------------------------
%  Codificación, idioma y paquetes
% ---------------------------------------------------------------------------
\usepackage[utf8]{inputenc}
\usepackage[T1]{fontenc}
\usepackage{lmodern}
\usepackage[spanish,es-noshorthands]{babel}
\usepackage{amsmath,amssymb}
\usepackage{array}
\usepackage{booktabs}
\usepackage{tikz}
\usetikzlibrary{arrows.meta,positioning,calc}

% ---------------------------------------------------------------------------
%  Tema y paleta (whale)
% ---------------------------------------------------------------------------
\usetheme{Boadilla}
\usecolortheme{whale}

\definecolor{upmblue}{RGB}{0,62,116}
\definecolor{upmcyan}{RGB}{0,132,180}
\definecolor{upmorange}{RGB}{214,93,14}
\definecolor{upmgray}{RGB}{90,98,108}
\definecolor{upmlight}{RGB}{235,241,247}

% "whale" deriva la paleta de structure.fg: al cambiarla, cambia todo el tema.
\setbeamercolor{structure}{fg=upmblue}
\setbeamercolor{alerted text}{fg=upmorange}

\setbeamercolor{title}{fg=white,bg=upmblue}
\setbeamercolor{subtitle}{fg=white,bg=upmblue}
\setbeamercolor{author}{fg=black!80}
\setbeamercolor{institute}{fg=black!60}
\setbeamercolor{date}{fg=black!60}

\setbeamercolor{block title}{fg=white,bg=upmblue}
\setbeamercolor{block body}{fg=black!88,bg=upmlight}
\setbeamercolor{block title alerted}{fg=white,bg=upmorange}
\setbeamercolor{block body alerted}{fg=black!88,bg=upmorange!8}
\setbeamercolor{block title example}{fg=white,bg=upmcyan}
\setbeamercolor{block body example}{fg=black!88,bg=upmcyan!8}

% Pie de página del outer theme "infolines" con tres tonos de azul.
\setbeamercolor{author in head/foot}{fg=white,bg=upmblue}
\setbeamercolor{title in head/foot}{fg=white,bg=upmblue!80!black}
\setbeamercolor{date in head/foot}{fg=white,bg=upmblue!60!black}

\setbeamercolor{itemize item}{fg=upmorange}
\setbeamercolor{itemize subitem}{fg=upmblue}
\setbeamercolor{itemize subsubitem}{fg=upmgray}

% Pie de página propio: tres bandas que suman el ancho de texto.
\setbeamertemplate{footline}{%
  \leavevmode%
  \hbox to \textwidth{%
    \begin{beamercolorbox}[wd=.38\textwidth,ht=2.4ex,dp=1.2ex,leftskip=.6em]{author in head/foot}%
      \usebeamerfont{author in head/foot}\insertshortauthor
    \end{beamercolorbox}%
    \begin{beamercolorbox}[wd=.44\textwidth,ht=2.4ex,dp=1.2ex,center]{title in head/foot}%
      \usebeamerfont{title in head/foot}\insertshorttitle
    \end{beamercolorbox}%
    \begin{beamercolorbox}[wd=.18\textwidth,ht=2.4ex,dp=1.2ex,rightskip=.6em]{date in head/foot}%
      \hfill\usebeamerfont{date in head/foot}\insertframenumber\,/\,\inserttotalframenumber
    \end{beamercolorbox}%
  }%
}

% ---------------------------------------------------------------------------
%  Tipografía y plantillas
% ---------------------------------------------------------------------------
\usefonttheme[onlymath]{serif}
\setbeamerfont{title}{series=\bfseries,size=\LARGE}
\setbeamerfont{subtitle}{size=\normalsize}
\setbeamerfont{frametitle}{series=\bfseries,size=\large}
\setbeamerfont{section title}{series=\bfseries,size=\Huge}

\setbeamertemplate{navigation symbols}{}
\setbeamertemplate{blocks}[rounded][shadow=false]
\setbeamertemplate{itemize item}{\small$\blacktriangleright$}
\setbeamertemplate{itemize subitem}{\scriptsize$\bullet$}

\setbeamersize{text margin left=1em,text margin right=1em}

% Portada de sección a sangre, en el azul de "whale".
\setbeamertemplate{section page}{%
  \begin{tikzpicture}[remember picture,overlay]
    \fill[upmblue] (current page.south west) rectangle (current page.north east);
    \node[anchor=west,text width=0.82\paperwidth,align=left]
      at ([xshift=0.09\paperwidth]current page.west) {%
        {\color{white!55}\fontsize{64}{64}\selectfont\bfseries\insertsectionnumber}\\[0.15cm]
        {\color{white}\fontsize{26}{30}\selectfont\bfseries\insertsectionhead}\\[0.35cm]
        {\color{upmorange}\rule{4.5cm}{3pt}}%
      };
  \end{tikzpicture}%
}
\AtBeginSection[]{%
  \begin{frame}[plain]
    \sectionpage
  \end{frame}%
}

% Cajas para el diagrama de flujo de la sección de híbridos.
\tikzset{
  flow/.style={
    draw=upmblue!70, fill=upmblue!8, rounded corners=3pt,
    align=center, minimum height=1.15cm, text width=3.35cm,
    font=\small, inner sep=4pt
  },
  flow arrow/.style={-{Latex[length=2.2mm]}, draw=upmorange, line width=1pt}
}

% Nodos y aristas de los grafos de ejemplo (introducción y algoritmo PC).
\tikzset{
  bn small/.style={
    circle, draw=upmblue, fill=upmlight, line width=0.8pt,
    font=\small, minimum size=0.62cm, inner sep=0pt
  },
  bn edge/.style={-{Latex[length=2mm]}, draw=upmblue, line width=0.9pt},
  skel/.style={draw=upmblue, line width=0.9pt},
  removed/.style={draw=upmorange, dashed, line width=0.9pt},
  limit/.style={
    draw=upmblue!70, fill=upmblue!8, rounded corners=3pt,
    align=center, text width=2.15cm, minimum height=1.9cm,
    font=\footnotesize, inner sep=3pt
  }
}

% ---------------------------------------------------------------------------
%  Macros de plantilla
% ---------------------------------------------------------------------------
% Texto de relleno: se sustituye por el contenido real.
\newcommand{\slot}[1]{\textcolor{upmgray}{\textit{#1}}}

% Hueco para una figura: \figslot[ancho]{texto}
\newcommand{\figslot}[2][0.9\linewidth]{%
  \begin{tikzpicture}
    \node[draw=upmgray!45, dashed, line width=0.8pt, rounded corners=4pt,
          minimum width=#1, minimum height=3.6cm, align=center,
          text=upmgray, font=\small] {#2};
  \end{tikzpicture}%
}

% ---------------------------------------------------------------------------
%  Metadatos
% ---------------------------------------------------------------------------
\title[Constraint-based learning]{Constraint-based learning}
\subtitle{Variantes del algoritmo PC, contrastes de independencia condicional\\y enfoques híbridos}
\author[D. Benasco, T. Camargo, J. D. Cárdenas]{%
  Damián Benasco Sánchez \and Thales Camargo \and Juan David Cárdenas}
\institute{Universidad Politécnica de Madrid}
\date{Máster Universitario en Inteligencia Artificial}

\graphicspath{{../report/Figures/}{./}}

\begin{document}

% ===========================================================================
%  PORTADA
% ===========================================================================
\begin{frame}[plain]
  \titlepage
\end{frame}

% ===========================================================================
%  INTRODUCCIÓN (antes del índice; no cuenta como sección)
% ===========================================================================

% ---------------------------------------------------------------------------
\begin{frame}[t]{Aprender la estructura: tres familias de métodos}
  \begin{columns}[T,onlytextwidth]
    \begin{column}{0.31\textwidth}
      \begin{block}{Puntuación y búsqueda}
        Evalúa grafos completos con una puntuación (p.\,ej.\ BIC) y busca el
        mejor.\\[0.2cm]
        {\small\color{upmgray}\itshape Como probarse ropa hasta encontrar la
        que mejor queda.}
      \end{block}
    \end{column}
    \begin{column}{0.31\textwidth}
      \begin{alertblock}{Basados en restricciones}
        Pruebas locales de independencia condicional; cada decisión es
        rastreable.\\[0.2cm]
        {\small\color{upmgray}\itshape Como un detective que descarta
        sospechosos con pistas.}
      \end{alertblock}
    \end{column}
    \begin{column}{0.31\textwidth}
      \begin{alertblock}{Híbridos}
        Las restricciones reducen el espacio y una puntuación refina el grafo
        (p.\,ej.\ MMHC).\\[0.2cm]
        {\small\color{upmgray}\itshape El detective hace la lista corta; el
        jurado decide.}
      \end{alertblock}
    \end{column}
  \end{columns}
  \vspace{0.4cm}
  \begin{exampleblock}{Tema del trabajo}
    De las restricciones a los híbridos: el PC, sus pruebas de independencia
    y su evolución hacia los métodos híbridos.
  \end{exampleblock}
\end{frame}

% ---------------------------------------------------------------------------
\begin{frame}[t]{Independencia condicional}
  \framesubtitle{¿$X$ e $Y$ son independientes si ya conozco $Z$?}
  \vspace{0.2cm}
  \begin{columns}[T,onlytextwidth]
    \begin{column}{0.31\textwidth}
      \centering\textbf{\color{upmblue}Cadena}\\[0.2cm]
      \begin{tikzpicture}[x=1.15cm]
        \node[bn small] (X) at (0,0) {$X$};
        \node[bn small] (Z) at (1,0) {$Z$};
        \node[bn small] (Y) at (2,0) {$Y$};
        \draw[bn edge] (X) -- (Z);
        \draw[bn edge] (Z) -- (Y);
      \end{tikzpicture}\\[0.2cm]
      $X\perp Y\mid Z$
    \end{column}
    \begin{column}{0.31\textwidth}
      \centering\textbf{\color{upmblue}Bifurcación}\\[0.2cm]
      \begin{tikzpicture}[x=1.15cm]
        \node[bn small] (X) at (0,0) {$X$};
        \node[bn small] (Z) at (1,0) {$Z$};
        \node[bn small] (Y) at (2,0) {$Y$};
        \draw[bn edge] (Z) -- (X);
        \draw[bn edge] (Z) -- (Y);
      \end{tikzpicture}\\[0.2cm]
      $X\perp Y\mid Z$
    \end{column}
    \begin{column}{0.31\textwidth}
      \centering\textbf{\color{upmorange}Colisor}\\[0.2cm]
      \begin{tikzpicture}[x=1.15cm]
        \node[bn small] (X) at (0,0) {$X$};
        \node[bn small] (Z) at (1,0) {$Z$};
        \node[bn small] (Y) at (2,0) {$Y$};
        \draw[bn edge] (X) -- (Z);
        \draw[bn edge] (Y) -- (Z);
      \end{tikzpicture}\\[0.2cm]
      $X\perp Y$, pero dependientes dado $Z$
    \end{column}
  \end{columns}
  \vspace{0.4cm}
  \begin{columns}[T,onlytextwidth]
    \begin{column}{0.48\textwidth}
      \begin{block}{Condición de Markov}
        El grafo implica independencias en los datos.
      \end{block}
    \end{column}
    \begin{column}{0.49\textwidth}
      \begin{block}{Fidelidad}
        Los datos no tienen independencias extra por coincidencias.
      \end{block}
    \end{column}
  \end{columns}
  \vfill
  {\scriptsize\color{upmgray}Los colisores dejan una huella distinta en los
  datos: son la clave para orientar las flechas.}
\end{frame}

% ===========================================================================
%  CONTENIDO
% ===========================================================================
\begin{frame}[t]{Contenido}
  \tableofcontents[hideallsubsections]
  \vfill
  \begin{block}{Hilo conductor}
    Aprender el grafo de una red bayesiana a partir de datos haciendo
    preguntas del tipo \emph{¿son $X$ e $Y$ independientes si ya conozco $Z$?}:
    el algoritmo PC, las pruebas estadísticas que responden a esas preguntas
    y los híbridos que corrigen su fragilidad.
  \end{block}
\end{frame}

% ===========================================================================
\section[El algoritmo PC]{El algoritmo PC: fundamentos, limitaciones y variantes}
% ===========================================================================

% ---------------------------------------------------------------------------
\begin{frame}[t]{Algoritmo PC: fundamentos}
  \framesubtitle{Peter Spirtes y Clark Glymour (1991): dos fases a partir de los datos.}
  \begin{center}
    \begin{tikzpicture}[node distance=0.55cm]
      \node[flow,text width=2.6cm] (datos) {\textbf{Datos}\\[2pt]
        {\footnotesize Observaciones de las variables}};
      \node[flow,text width=2.6cm,right=of datos] (f1) {\textbf{Fase 1}\\[2pt]
        {\footnotesize Esqueleto: qué está conectado}};
      \node[flow,text width=2.6cm,right=of f1] (f2) {\textbf{Fase 2}\\[2pt]
        {\footnotesize Orientación: hacia dónde van las flechas}};
      \node[flow,text width=2.6cm,right=of f2,fill=upmorange!10,draw=upmorange]
        (cpdag) {\textbf{CPDAG}\\[2pt]
        {\footnotesize Grafo parcialmente dirigido}};
      \draw[flow arrow] (datos) -- (f1);
      \draw[flow arrow] (f1) -- (f2);
      \draw[flow arrow] (f2) -- (cpdag);
    \end{tikzpicture}
  \end{center}
  \begin{itemize}
    \item Supuestos: condición de Markov causal, fidelidad y suficiencia
          causal (sin variables ocultas).
    \item Fase 1: parte del grafo completo y poda aristas con contrastes de
          independencia, guardando el conjunto separador de cada par.
    \item Fase 2: orienta $v$-estructuras (colisores) y propaga con las reglas
          de Meek.
  \end{itemize}
  \begin{alertblock}{La idea que lo hace eficiente}
    Solo condiciona sobre los vecinos actuales de cada variable, no sobre
    todos los subconjuntos posibles (como hacía SGS). Coste exponencial en el
    peor caso, pero polinómico si el grafo es poco denso.
  \end{alertblock}
\end{frame}

% ---------------------------------------------------------------------------
\begin{frame}[t]{Fase 1: el esqueleto}
  \framesubtitle{Todo conectado al inicio; se eliminan aristas con pruebas.}
  {\small\color{upmgray}Grafo verdadero (desconocido): $A\to C\leftarrow B$,
  $C\to D$.}
  \vspace{0.2cm}
  \begin{columns}[T,onlytextwidth]
    \begin{column}{0.31\textwidth}
      \centering\textbf{\color{upmblue}1. Grafo completo}\\[0.25cm]
      \begin{tikzpicture}[x=1.1cm,y=0.9cm]
        \node[bn small] (A) at (0,2) {$A$};
        \node[bn small] (B) at (2,2) {$B$};
        \node[bn small] (C) at (1,1) {$C$};
        \node[bn small] (D) at (1,0) {$D$};
        \draw[skel] (A) -- (B);
        \draw[skel] (A) -- (C);
        \draw[skel] (B) -- (C);
        \draw[skel] (C) -- (D);
        \draw[skel] (A) to[bend right=25] (D);
        \draw[skel] (B) to[bend left=25] (D);
      \end{tikzpicture}\\[0.25cm]
      {\small Todo conectado.}
    \end{column}
    \begin{column}{0.31\textwidth}
      \centering\textbf{\color{upmblue}2. Orden $\ell=0$}\\[0.25cm]
      \begin{tikzpicture}[x=1.1cm,y=0.9cm]
        \node[bn small] (A) at (0,2) {$A$};
        \node[bn small] (B) at (2,2) {$B$};
        \node[bn small] (C) at (1,1) {$C$};
        \node[bn small] (D) at (1,0) {$D$};
        \draw[removed] (A) -- (B);
        \draw[skel] (A) -- (C);
        \draw[skel] (B) -- (C);
        \draw[skel] (C) -- (D);
        \draw[skel] (A) to[bend right=25] (D);
        \draw[skel] (B) to[bend left=25] (D);
      \end{tikzpicture}\\[0.25cm]
      {\small $A\perp B$ $\Rightarrow$ se elimina $A$--$B$\\
      $S(A,B)=\emptyset$}
    \end{column}
    \begin{column}{0.31\textwidth}
      \centering\textbf{\color{upmblue}3. Orden $\ell=1$}\\[0.25cm]
      \begin{tikzpicture}[x=1.1cm,y=0.9cm]
        \node[bn small] (A) at (0,2) {$A$};
        \node[bn small] (B) at (2,2) {$B$};
        \node[bn small] (C) at (1,1) {$C$};
        \node[bn small] (D) at (1,0) {$D$};
        \draw[skel] (A) -- (C);
        \draw[skel] (B) -- (C);
        \draw[skel] (C) -- (D);
        \draw[removed] (A) to[bend right=25] (D);
        \draw[removed] (B) to[bend left=25] (D);
      \end{tikzpicture}\\[0.25cm]
      {\small $A\perp D\mid C$, \ $B\perp D\mid C$\\
      $S(A,D)=S(B,D)=\{C\}$}
    \end{column}
  \end{columns}
  \vfill
  {\scriptsize\color{upmgray}Líneas discontinuas: aristas eliminadas en cada
  paso. Se guarda el conjunto separador de cada par, porque la fase 2 lo
  usa.}
\end{frame}

% ---------------------------------------------------------------------------
\begin{frame}[t]{Fase 2: orientar las flechas}
  \framesubtitle{Primero los colisores, después reglas de propagación.}
  \vspace{0.3cm}
  \begin{columns}[T,onlytextwidth]
    \begin{column}{0.30\textwidth}
      \centering\textbf{\color{upmblue}1. Colisor}\\[0.25cm]
      \begin{tikzpicture}[x=1.1cm,y=0.9cm]
        \node[bn small] (A) at (0,2) {$A$};
        \node[bn small] (B) at (2,2) {$B$};
        \node[bn small] (C) at (1,1) {$C$};
        \node[bn small] (D) at (1,0) {$D$};
        \draw[bn edge,upmorange] (A) -- (C);
        \draw[bn edge,upmorange] (B) -- (C);
        \draw[skel] (C) -- (D);
      \end{tikzpicture}\\[0.25cm]
      {\small $C\notin S(A,B)=\emptyset$\\
      $\Rightarrow A\to C\leftarrow B$}
    \end{column}
    \begin{column}{0.30\textwidth}
      \centering\textbf{\color{upmblue}2. Propagación (Meek)}\\[0.25cm]
      \begin{tikzpicture}[x=1.1cm,y=0.9cm]
        \node[bn small] (A) at (0,2) {$A$};
        \node[bn small] (B) at (2,2) {$B$};
        \node[bn small] (C) at (1,1) {$C$};
        \node[bn small] (D) at (1,0) {$D$};
        \draw[bn edge] (A) -- (C);
        \draw[bn edge] (B) -- (C);
        \draw[bn edge,upmorange] (C) -- (D);
      \end{tikzpicture}\\[0.25cm]
      {\small Sin ciclos ni colisores nuevos\\
      $\Rightarrow C\to D$}
    \end{column}
    \begin{column}{0.34\textwidth}
      \begin{block}{Resultado: un CPDAG}
        Las aristas que los datos no permiten orientar quedan sin dirección:
        \[
          X\to Y\to Z \quad\text{y}\quad X\leftarrow Y\leftarrow Z
        \]
        implican las mismas independencias.
      \end{block}
    \end{column}
  \end{columns}
  \vfill
  {\scriptsize\color{upmgray}Intuición: si $C$ fuera intermediario o causa
  común, habría hecho falta $C$ para separar $A$ y $B$.}
\end{frame}

% ---------------------------------------------------------------------------
\begin{frame}[t]{Del oráculo a los datos}
  \framesubtitle{Con pruebas perfectas acierta; con datos reales puede fallar.}
  \begin{columns}[T,onlytextwidth]
    \begin{column}{0.48\textwidth}
      \begin{exampleblock}{Oráculo ideal}
        Pruebas perfectas: el PC es correcto y completo.
      \end{exampleblock}
    \end{column}
    \begin{column}{0.49\textwidth}
      \begin{alertblock}{Datos reales}
        Cada prueba puede fallar y el error se propaga a las siguientes.
      \end{alertblock}
    \end{column}
  \end{columns}
  \vspace{0.3cm}
  \begin{center}
    \begin{tikzpicture}[node distance=0.25cm]
      \node[limit] (l1) {\textbf{\color{upmorange}(i)}\\[1pt]\textbf{Alta dimensión}\\[1pt]
        {\scriptsize muchas variables, pocos datos}};
      \node[limit,right=of l1] (l2) {\textbf{\color{upmorange}(ii)}\\[1pt]\textbf{Orden de las variables}\\[1pt]
        {\scriptsize el resultado cambia al reordenar}};
      \node[limit,right=of l2] (l3) {\textbf{\color{upmorange}(iii)}\\[1pt]\textbf{Conjuntos grandes}\\[1pt]
        {\scriptsize pruebas poco fiables al condicionar mucho}};
      \node[limit,right=of l3] (l4) {\textbf{\color{upmorange}(iv)}\\[1pt]\textbf{Coste computacional}\\[1pt]
        {\scriptsize miles de pruebas}};
      \node[limit,right=of l4] (l5) {\textbf{\color{upmorange}(v)}\\[1pt]\textbf{Supuestos rígidos}\\[1pt]
        {\scriptsize i.i.d., fidelidad, sin variables ocultas}};
    \end{tikzpicture}
  \end{center}
  \vfill
  {\scriptsize\color{upmgray}Como el algoritmo es secuencial, un error al
  principio cambia los vecinos y, con ellos, las pruebas posteriores. Estas
  cinco limitaciones organizan el resto de la sección.}
\end{frame}

% ---------------------------------------------------------------------------
\begin{frame}[t]{(i) Alta dimensión}
  \framesubtitle{¿Funciona el PC con más variables que observaciones?}
  \vspace{0.2cm}
  \begin{columns}[T,onlytextwidth]
    \begin{column}{0.30\textwidth}
      \centering
      {\color{upmblue}\fontsize{40}{44}\selectfont\bfseries $p\gg n$}\\[0.3cm]
      {\small más variables ($p$) que observaciones ($n$)}
    \end{column}
    \begin{column}{0.66\textwidth}
      \begin{block}{Kalisch y Bühlmann (2007)}
        El PC es consistente en alta dimensión si:
        \begin{itemize}
          \item el grafo es poco denso (pocos vecinos);
          \item los datos son gaussianos;
          \item el nivel de significación $\alpha$ disminuye con $n$.
        \end{itemize}
      \end{block}
      \begin{alertblock}{La letra pequeña: fidelidad fuerte \textperiodcentered{} Uhler et al.\ (2013)}
        La prueba exige que las dependencias reales no sean demasiado débiles.
        Las distribuciones que incumplen esta condición no son casos raros:
        pueden ser una fracción sorprendentemente grande. Por eso el PC es
        frágil con muestras finitas.
      \end{alertblock}
    \end{column}
  \end{columns}
\end{frame}

% ---------------------------------------------------------------------------
\begin{frame}[t]{(ii) PC-stable y dependencia del orden}
  \framesubtitle{Reordenar las columnas no debería cambiar el resultado.}
  \begin{columns}[T,onlytextwidth]
    \begin{column}{0.48\textwidth}
      \begin{alertblock}{El problema: PC original}
        \begin{itemize}
          \item Elimina una arista y actualiza los vecinos al instante.
          \item Las pruebas siguientes dependen de lo que ya se eliminó.
          \item Distinto orden de variables $\Rightarrow$ grafo distinto.
        \end{itemize}
      \end{alertblock}
    \end{column}
    \begin{column}{0.49\textwidth}
      \begin{block}{La solución: PC-stable \textperiodcentered{} Colombo y Maathuis (2014)}
        \begin{itemize}
          \item Congela la lista de vecinos al inicio de cada nivel.
          \item Todas las pruebas del nivel usan la misma lista.
          \item Esqueleto independiente del orden.
        \end{itemize}
      \end{block}
    \end{column}
  \end{columns}
  \vspace{0.2cm}
  \begin{exampleblock}{Para las flechas: CPC (Ramsey et al., 2006)}
    Orienta un colisor solo si todos los separadores coinciden; si no, lo
    marca como ambiguo.
  \end{exampleblock}
  \vfill
  {\small\color{upmorange}\textbf{Precio:} más pruebas
  $\;\longrightarrow\;$ (iv)}
\end{frame}

% ---------------------------------------------------------------------------
\begin{frame}[t]{(iii) Conjuntos de condicionamiento grandes}
  \framesubtitle{Condicionar en muchas variables deja pocos datos por celda.}
  \begin{columns}[T,onlytextwidth]
    \begin{column}{0.48\textwidth}
      \centering
      \begin{tikzpicture}[x=0.75cm,y=0.0135cm]
        \foreach \s/\v/\lab in {0/250/250,1/125/125,2/62.5/62,3/31.25/31,
                                4/15.6/16,5/7.8/8,6/3.9/3.9} {
          \fill[upmblue!75] (\s-0.32,0) rectangle (\s+0.32,\v);
          \node[above,font=\scriptsize] at (\s,\v) {\lab};
          \node[below,font=\scriptsize] at (\s,0) {\s};
        }
        \fill[upmorange] (6-0.32,0) rectangle (6+0.32,3.9);
        \draw[upmgray] (-0.6,0) -- (6.6,0);
        \node[font=\scriptsize,text=upmgray] at (3,-28) {$|S|$};
        \node[font=\scriptsize,text=upmgray,anchor=west] at (1.6,230)
          {datos por celda};
        \node[font=\scriptsize,text=upmgray,anchor=west] at (1.6,205)
          {($n=1000$, binarias)};
      \end{tikzpicture}\\[0.1cm]
      {\small Con $|S|=6$ ya hay 256 celdas: menos de 4 datos en cada una.}
    \end{column}
    \begin{column}{0.49\textwidth}
      \begin{block}{Abellán et al.\ (2006)}
        Cortes mínimos para usar conjuntos más pequeños y una prueba bayesiana
        (BD) en lugar de $\chi^2$.
      \end{block}
      \begin{block}{$k$-PC \textperiodcentered{} Kocaoglu (2023)}
        Solo pruebas con $|S|\le k$: aprende una clase algo menos precisa,
        pero comete menos errores.
      \end{block}
      {\small\color{upmgray}\itshape Aprender menos, pero aprenderlo bien.}
    \end{column}
  \end{columns}
\end{frame}

% ---------------------------------------------------------------------------
\begin{frame}[t]{(iv) Coste computacional}
  \framesubtitle{Las pruebas de un nivel de PC-stable son independientes: paralelizables.}
  \vspace{0.2cm}
  \begin{columns}[T,onlytextwidth]
    \begin{column}{0.44\textwidth}
      \centering
      {\color{upmblue}\fontsize{22}{26}\selectfont\bfseries
      ${>}\,24$\,h $\to$ {\color{upmorange}${<}\,6$\,h}}\\[0.2cm]
      {\small Parallel-PC (Le et al., 2019)\\
      DREAM5: más de 1600 variables, 8 núcleos}\\[0.4cm]
      \begin{tikzpicture}[x=1.25cm,y=0.42cm]
        \foreach \c in {1,...,4} {
          \node[font=\scriptsize\bfseries,text=upmblue] at (\c,0.9) {Núcleo \c};
          \foreach \t in {0,1,2} {
            \node[draw=upmblue!70,fill=upmblue!8,rounded corners=2pt,
                  font=\tiny,minimum width=1.05cm,minimum height=0.32cm]
              at (\c,-\t) {test};
          }
        }
      \end{tikzpicture}
    \end{column}
    \begin{column}{0.52\textwidth}
      \begin{block}{Madsen et al.\ (2017)}
        Hilos con bloques de variables y procesos MPI que se reparten los
        datos.
      \end{block}
      \begin{block}{Scutari (2017) \textperiodcentered{} \texttt{bnlearn}}
        El \emph{backtracking} empeora la estabilidad; paralelizar es mejor.
      \end{block}
      \begin{block}{Le et al.\ (2019)}
        Reparte entre núcleos las pruebas de cada nivel de PC-stable, sin
        perder la independencia del orden.
      \end{block}
    \end{column}
  \end{columns}
\end{frame}

% ---------------------------------------------------------------------------
\begin{frame}[t]{(v) Relajar los supuestos}
  \framesubtitle{Llevar el PC a datos que no cumplen sus hipótesis.}
  \vspace{0.1cm}
  \begin{columns}[T,onlytextwidth]
    \begin{column}{0.31\textwidth}
      \begin{block}{1. Series temporales}
        {\small\textbf{TPC} \textperiodcentered{} Biswas y Mukherjee (2024)\\
        \textbf{PCMCI} \textperiodcentered{} Runge et al.\ (2019)}\\[0.15cm]
        Datos que dependen del tiempo; relaciones con retardo (la lluvia de
        ayer afecta al caudal de hoy).
      \end{block}
    \end{column}
    \begin{column}{0.31\textwidth}
      \begin{block}{2. Eventos extremos}
        {\small\textbf{PCstar} \textperiodcentered{} Améndola et al.\ (2026)}\\[0.15cm]
        Modelos de colas pesadas (riesgo financiero, crecidas) donde la
        fidelidad a la $d$-separación falla.
      \end{block}
    \end{column}
    \begin{column}{0.31\textwidth}
      \begin{block}{3. Decisiones suaves}
        {\small\textbf{DAGPA} \textperiodcentered{} Zhou et al.\ (2025)}\\[0.15cm]
        Convierte la $d$-separación en una función continua y optimiza el
        grafo con gradientes.
      \end{block}
    \end{column}
  \end{columns}
  \vfill
  {\scriptsize\color{upmgray}Variables ocultas: FCI y RFCI (se tratan en la
  sección de pruebas de independencia).}
\end{frame}

% ---------------------------------------------------------------------------
\begin{frame}[t]{Evolución del algoritmo PC}
  \footnotesize
  \renewcommand{\arraystretch}{1.12}
  \begin{tabular}{@{}p{0.10\textwidth}p{0.26\textwidth}p{0.35\textwidth}p{0.25\textwidth}@{}}
    \toprule
    \textbf{Año} & \textbf{Variante} & \textbf{Idea clave} & \textbf{Referencia} \\
    \midrule
    1991    & PC                   & Esqueleto + orientación              & Spirtes y Glymour \\
    1995    & Reglas de orientación & Propagación sin ciclos               & Meek \\
    2006    & CPC, cortes mínimos  & Fiabilidad estadística (ii, iii)     & Ramsey; Abellán \\
    2007    & Alta dimensión       & Consistencia con $p\gg n$ (i)        & Kalisch y Bühlmann \\
    2014    & PC-stable            & Independencia del orden (ii)         & Colombo y Maathuis \\
    2017--19 & Paralelización      & Pruebas de un nivel en paralelo (iv) & Madsen; Scutari; Le \\
    2023    & $k$-PC               & Conjuntos pequeños (iii)             & Kocaoglu \\
    2024--26 & Tiempo, extremos, DAGPA & Relajar supuestos (v)            & Biswas; Améndola; Zhou \\
    \bottomrule
  \end{tabular}
  \vspace{0.3cm}
  \begin{block}{Criterio de agrupación}
    Cada mejora responde a una limitación concreta: el PC no se ha
    reemplazado, se ha ido parcheando problema a problema.
  \end{block}
\end{frame}

% ---------------------------------------------------------------------------
\begin{frame}[t]{Resumen del algoritmo PC}
  \framesubtitle{Tres ideas antes de pasar a las pruebas de independencia.}
  \vspace{0.2cm}
  \begin{columns}[T,onlytextwidth]
    \begin{column}{0.31\textwidth}
      \begin{block}{1. Interpretable}
        El PC deduce el grafo con pruebas de independencia; devuelve un CPDAG,
        no un único DAG.
      \end{block}
    \end{column}
    \begin{column}{0.31\textwidth}
      \begin{alertblock}{2. Frágil}
        Los errores estadísticos se propagan; cada variante corrige una cara
        del problema.
      \end{alertblock}
    \end{column}
    \begin{column}{0.31\textwidth}
      \begin{exampleblock}{3. Depende de sus pruebas}
        Cada decisión es una prueba estadística: elegirla bien es clave
        (sección~2).
      \end{exampleblock}
    \end{column}
  \end{columns}
  \vfill
  {\scriptsize\color{upmgray}A menudo resolver un problema genera otro, como
  PC-stable, que hace más pruebas. Esta fragilidad motiva también los
  métodos híbridos.}
\end{frame}

% ===========================================================================
\section[Pruebas de independencia condicional]{Pruebas de independencia condicional y extensiones para causalidad}

% ===========================================================================
%  2.2  INTRODUCCIÓN
% ===========================================================================

\begin{frame}[t]{Pruebas de IC: el motor estadístico}
  \begin{columns}[T,onlytextwidth]
    \begin{column}{0.53\textwidth}
      \begin{itemize}
        \item Cada decisión de poda u orientación del PC se apoya en
              contrastes $H_0: I(X,Y\mid\mathbf Z)$.
        \item Bajo Markov causal y fidelidad:
      \end{itemize}
      \[
        I(X,Y\mid\mathbf Z)
        \;\Longleftrightarrow\;
        X\perp_G Y\mid\mathbf Z.
      \]
      \begin{itemize}
        \item Frente a las funciones de puntuación, que evalúan estructuras
              completas, la decisión es \alert{local y trazable}: cada arista
              se justifica con uno o varios contrastes y sus $p$-valores.
        \item No hay contraste universal: van de los paramétricos clásicos a
              los no paramétricos, y cada uno impone un modelo de dependencia
              y paga esa restricción con potencia.
      \end{itemize}
    \end{column}
    \begin{column}{0.44\textwidth}
      \begin{block}{Por qué importa}
        La calidad del grafo aprendido queda acotada por la calidad de las
        decisiones estadísticas.
      \end{block}
      \begin{alertblock}{Un contraejemplo}
        Si $X\sim\mathcal N(0,1)$ e $Y=X^2$:
        \[
          \mathrm{corr}(X,Y)=0
          \quad\text{pero}\quad
          \neg I(X,Y).
        \]
        La correlación parcial no la detecta.
      \end{alertblock}
      {\scriptsize\color{upmgray}Cuatro bloques de contrastes: parcial,
      tablas, regresión/GCM y no paramétricos (kernel y kNN).}
    \end{column}
  \end{columns}
\end{frame}

% ===========================================================================
%  2.2.1  CORRELACIÓN PARCIAL
% ===========================================================================
\subsection{Correlación parcial}

\begin{frame}[t]{Correlación parcial: el contraste gaussiano}
  \begin{columns}[T,onlytextwidth]
    \begin{column}{0.55\textwidth}
      \begin{itemize}
        \item Datos continuos con normalidad conjunta.
        \item $I(X,Y\mid\mathbf Z)$ equivale a que la correlación parcial
              $\rho_{XY\cdot\mathbf Z}$ se anule.
        \item Se estima la covarianza muestral $S$ de $(X,Y,\mathbf Z)$ y se
              invierte: $\hat K=S^{-1}$.
      \end{itemize}
      \[
        \hat\rho_{XY\cdot\mathbf Z}
        =\frac{-\hat K_{XY}}{\sqrt{\hat K_{XX}\hat K_{YY}}}.
      \]
      \begin{itemize}
        \item La transformada $z$ de Fisher convierte la correlación en un
              estadístico aproximadamente normal.
      \end{itemize}
    \end{column}
    \begin{column}{0.42\textwidth}
      \begin{block}{Contraste}
        \[
          z(r)=\tfrac12\ln\frac{1+r}{1-r},
        \]
        \[
          T=\sqrt{n-|\mathbf Z|-3}\;z(\hat\rho).
        \]
        Se rechaza $H_0$ si $|T|>z_{1-\alpha/2}$.
      \end{block}
      \begin{exampleblock}{Coste}
        $\mathcal O(|\mathbf Z|^3)$ por contraste, con $p$-valores
        analíticos.
      \end{exampleblock}
    \end{column}
  \end{columns}
\end{frame}

\begin{frame}[t]{Correlación parcial: ventajas y limitaciones}
  \begin{columns}[T,onlytextwidth]
    \begin{column}{0.48\textwidth}
      \begin{block}{Ventajas}
        \begin{itemize}
          \item Coste bajo: invertir el bloque de covarianzas.
          \item $p$-valores analíticos, sin remuestreo.
          \item Implementación de referencia en \texttt{pcalg} y
                \texttt{bnlearn}.
          \item Sostiene la consistencia del PC en alta dimensión con grafos
                dispersos.
        \end{itemize}
      \end{block}
    \end{column}
    \begin{column}{0.49\textwidth}
      \begin{alertblock}{Limitaciones}
        \begin{itemize}
          \item Solo es válido bajo normalidad y linealidad.
          \item Fuera de ese modelo, las versiones de rango solo aproximan.
          \item Correlación nula no implica independencia en general.
          \item En muestras pequeñas, la aproximación asintótica del
                $p$-valor puede ser pobre.
        \end{itemize}
      \end{alertblock}
    \end{column}
  \end{columns}
  \vspace{0.2cm}
  \begin{exampleblock}{Líneas recientes}
    WLS para varianza no constante; correlaciones de rango y cópulas para
    no normalidad; permutaciones para calibrar en muestra pequeña.
  \end{exampleblock}
\end{frame}

% ===========================================================================
%  2.2.2  TABLAS DE CONTINGENCIA
% ===========================================================================
\subsection{Tablas de contingencia}

\begin{frame}[t]{Tablas de contingencia: el contraste discreto}
  \begin{columns}[T,onlytextwidth]
    \begin{column}{0.57\textwidth}
      \begin{itemize}
        \item Datos discretos o categóricos: se comparan los conteos observados
              con los esperados bajo independencia.
        \item $n_{ijk}$: observaciones con $X=i$, $Y=j$ y $\mathbf Z=k$.
      \end{itemize}
      \begin{block}{Estadísticos}
        \[
          E_{ijk}=\frac{n_{i\cdot k}\,n_{\cdot jk}}{n_{\cdot\cdot k}},
        \]
        \[
          \chi^2=\sum_{i,j,k}\frac{(n_{ijk}-E_{ijk})^2}{E_{ijk}},
        \]
        \[
          G^2=2\sum_{i,j,k} n_{ijk}\ln\frac{n_{ijk}}{E_{ijk}}.
        \]
      \end{block}
    \end{column}
    \begin{column}{0.40\textwidth}
      \begin{block}{Distribución nula}
        \[
          \nu=(I-1)(J-1)\prod_l K_l
        \]
        Se rechaza $H_0$ si el estadístico supera
        $\chi^2_{\nu,1-\alpha}$.
      \end{block}
      \begin{exampleblock}{Lectura}
        $G^2=2n\,\hat I(X,Y\mid\mathbf Z)$: versión discreta del contraste de
        razón de verosimilitudes, dentro de la familia de información mutua.
      \end{exampleblock}
      {\scriptsize\color{upmgray}Ordinales: tendencia de Jonckheere--Terpstra.}
    \end{column}
  \end{columns}
\end{frame}

\begin{frame}[t]{Tablas de contingencia: ventajas y limitaciones}
  \begin{columns}[T,onlytextwidth]
    \begin{column}{0.48\textwidth}
      \begin{block}{Ventajas}
        \begin{itemize}
          \item Contraste canónico para datos categóricos.
          \item $p$-valores analíticos.
          \item Lectura directa como divergencia entre la tabla observada y la
                esperada.
          \item En las simulaciones originales del PC, $G^2$ acertó el grafo
                más veces que $\chi^2$.
        \end{itemize}
      \end{block}
    \end{column}
    \begin{column}{0.49\textwidth}
      \begin{alertblock}{Limitaciones}
        \begin{itemize}
          \item Celdas $I\cdot J\cdot\prod_l K_l$: crecimiento exponencial con
                $|\mathbf Z|$ y la cardinalidad.
          \item Con pocas observaciones por celda, la aproximación $\chi^2$
                deja de ser fiable.
          \item Regla práctica: $n\ge 5\times$ número de celdas.
          \item El coste es lineal en celdas, pero las celdas son el cuello de
                botella.
        \end{itemize}
      \end{alertblock}
    \end{column}
  \end{columns}
  \vspace{0.1cm}
  \begin{exampleblock}{Líneas recientes}
    Tablas suavizadas (Dash et al., 2003); SECMI (Kubkowski et al., 2021);
    permutaciones.
  \end{exampleblock}
\end{frame}

% ===========================================================================
%  2.2.3  REGRESIÓN Y RESIDUOS
% ===========================================================================
\subsection{Regresión y métodos de residuos}

\begin{frame}[t]{Regresión y residuos: el contraste GCM}
  \begin{columns}[T,onlytextwidth]
    \begin{column}{0.53\textwidth}
      \begin{itemize}
        \item Datos mixtos: no se modela directamente la dependencia, sino que
              se elimina el efecto de $\mathbf Z$.
        \item La versión general es la \alert{medida de covarianza
              generalizada (GCM)} de Shah \& Peters (2020).
        \item Se ajustan $\hat m_X(\mathbf z)$ y $\hat m_Y(\mathbf z)$ y se
              toman los residuos
              $r_i=x_i-\hat m_X(\mathbf z_i)$ y
              $s_i=y_i-\hat m_Y(\mathbf z_i)$.
        \item Bajo independencia, los productos $R_i=r_i s_i$ tienen media
              cero.
      \end{itemize}
      \[
        T_n=\frac{\sqrt n\,\bar R}{\hat\sigma_R},
        \qquad
        \hat\sigma_R^2=\frac1n\sum_{i=1}^n(R_i-\bar R)^2.
      \]
    \end{column}
    \begin{column}{0.44\textwidth}
      \begin{block}{Distribución nula}
        $T_n\longrightarrow\mathcal N(0,1)$ si el producto de los errores
        cuadráticos medios de ambas regresiones es $o(n^{-1})$.
      \end{block}
      \begin{exampleblock}{Caso lineal-gaussiano}
        El contraste coincide, salvo normalización, con la correlación parcial
        de Fisher.
      \end{exampleblock}
      \begin{alertblock}{Idea}
        No fija la forma de la dependencia, sino la de las medias
        condicionales.
      \end{alertblock}
    \end{column}
  \end{columns}
\end{frame}

\begin{frame}[t]{Regresión y residuos: ventajas y limitaciones}
  \begin{columns}[T,onlytextwidth]
    \begin{column}{0.48\textwidth}
      \begin{block}{Ventajas}
        \begin{itemize}
          \item Flexible: lineal, GLM, bosques y kernels; admite datos mixtos.
          \item Cubre parte de la no linealidad.
          \item WGCM pondera por la varianza; el doble ML evita el sobreajuste.
        \end{itemize}
      \end{block}
    \end{column}
    \begin{column}{0.49\textwidth}
      \begin{alertblock}{Limitaciones}
        \begin{itemize}
          \item Las regresiones no son neutrales: solo detectan alternativas
                con covarianza residual.
          \item Sin restringir la clase nula, ningún contraste con nivel
                uniforme tiene potencia (Shah \& Peters, 2020).
          \item Elegir las regresiones exige conocer el problema; su coste
                puede ser alto.
        \end{itemize}
      \end{alertblock}
    \end{column}
  \end{columns}
  \begin{exampleblock}{Líneas recientes}
    WGCM amplía las alternativas a costa de potencia; el doble ML permite
    regresiones flexibles y calibradas.
  \end{exampleblock}
\end{frame}

% ===========================================================================
%  2.2.4  MÉTODOS NO PARAMÉTRICOS
% ===========================================================================
\subsection{Métodos no paramétricos}

\begin{frame}[t]{Kernel y vecinos: contrastes no paramétricos}
  \begin{columns}[T,onlytextwidth]
    \begin{column}{0.48\textwidth}
      \begin{block}{Kernel: HSIC y KCI}
        \begin{itemize}
          \item HSIC: dos variables son independientes si y solo si ninguna
                función de una covaría con ninguna función de la otra.
          \item KCI elimina $\mathbf Z$ por regresión no paramétrica y aplica
                una medida tipo HSIC a los residuos.
          \item Nulo por Monte Carlo o aproximación gamma.
        \end{itemize}
        \[
          \mathrm{HSIC}(X,Y)=0\iff X\perp Y.
        \]
      \end{block}
    \end{column}
    \begin{column}{0.49\textwidth}
      \begin{block}{Vecinos más cercanos}
        \begin{itemize}
          \item La información mutua condicional se estima con un estimador
                de $k$ vecinos.
          \item El contraste se calibra con permutaciones locales.
          \item RCIT/RCoT sustituyen los kernels por características de
                Fourier aleatorias y pasan a coste lineal.
        \end{itemize}
      \end{block}
    \end{column}
  \end{columns}
  \vspace{-2pt}
  {\small
  \begin{alertblock}{Idea común}
    No se fija la forma de la dependencia ni la de las medias condicionales.
  \end{alertblock}}
\end{frame}

\begin{frame}[t]{Métodos no paramétricos: ventajas y limitaciones}
  \begin{columns}[T,onlytextwidth]
    \begin{column}{0.48\textwidth}
      \begin{block}{Ventajas}
        \begin{itemize}
          \item No asumen forma funcional.
          \item Captan dependencias no lineales.
          \item Con densidades no suaves, los vecinos superan en potencia a
                los métodos de kernel.
          \item Los de kernel mitigan la maldición de la dimensionalidad;
                RCIT/RCoT escalan a muestras grandes.
        \end{itemize}
      \end{block}
    \end{column}
    \begin{column}{0.49\textwidth}
      \begin{alertblock}{Limitaciones}
        \begin{itemize}
          \item KCI es $\mathcal O(n^2)$ o peor: inviable en muestras grandes.
          \item La distribución nula analítica no está bien calibrada por
                debajo de $\sim 1000$ muestras.
          \item Las permutaciones mejoran la calibración a costa de potencia.
        \end{itemize}
      \end{alertblock}
    \end{column}
  \end{columns}
  \vspace{0.2cm}
  \begin{exampleblock}{Líneas recientes}
    Contrastes \emph{distribution-free} cuyo índice se anula si y solo si hay
    independencia; permutaciones aprendidas para estimar mejor el nulo.
  \end{exampleblock}
\end{frame}

% ===========================================================================
%  2.2.5  EXTENSIONES PARA CAUSALIDAD
% ===========================================================================
\subsection{Extensiones para causalidad}

\begin{frame}[t]{De la independencia a la causalidad}
  \begin{columns}[T,onlytextwidth]
    \begin{column}{0.47\textwidth}
      \begin{block}{El problema de encadenar contrastes}
        \begin{itemize}
          \item Ni $\alpha$ ni la potencia miden el error de la búsqueda, que
                depende del número de contrastes y de la estructura recorrida.
          \item Una eliminación errónea inhibe contrastes posteriores y
                arrastra esqueleto y orientación.
          \item La consistencia exige $\alpha_n\to 0$ con $n$.
        \end{itemize}
      \end{block}
    \end{column}
    \begin{column}{0.50\textwidth}
      \begin{block}{Nuevos criterios de separación}
        \begin{itemize}
          \item \textbf{$\sigma$-separación}: mecanismos no lineales y ciclos.
          \item \textbf{$m$-separación} en MAG/PAG: variables latentes y sesgo
                de selección (FCI).
        \end{itemize}
      \end{block}
      \begin{block}{Nuevos escenarios}
        \begin{itemize}
          \item \textbf{COmbINE}: unifica datos observacionales e
                intervencionales parcialmente solapados.
          \item \textbf{CTBN}: contrastes para trayectorias continuas; explotan
                la asimetría del tiempo.
        \end{itemize}
      \end{block}
    \end{column}
  \end{columns}
  \vspace{0.1cm}
  {\scriptsize\color{upmgray}Forré \& Mooij (2018); Richardson \& Spirtes
  (2002); Triantafillou \& Tsamardinos (2015); Villa \& Stella (2016).}
\end{frame}

\begin{frame}[t]{Extensiones causales: alcance y límites}
  \begin{columns}[T,onlytextwidth]
    \begin{column}{0.48\textwidth}
      \begin{block}{Ventajas}
        \begin{itemize}
          \item Dotan a los contrastes de semántica causal plena.
          \item Cubren ciclos, latentes y sesgo de selección.
          \item Integran datos heterogéneos y temporales.
          \item FCI/RFCI y COmbINE amplían el marco sin abandonar las
                restricciones.
        \end{itemize}
      \end{block}
    \end{column}
    \begin{column}{0.49\textwidth}
      \begin{alertblock}{Limitaciones}
        \begin{itemize}
          \item La vulnerabilidad al encadenamiento en muestra finita
                permanece intacta.
          \item Un solo error estadístico altera los conjuntos de
                condicionamiento posteriores.
          \item Las garantías son asintóticas.
          \item Cada extensión añade supuestos y requisitos de datos.
        \end{itemize}
      \end{alertblock}
    \end{column}
  \end{columns}
  \vspace{0.2cm}
  \begin{exampleblock}{Cierre de la sección}
    Esta fragilidad estructural es la motivación central de los enfoques
    híbridos de la sección~3.
  \end{exampleblock}
\end{frame}

% ============================================================================
% SOLO SECCION 3 - Pegar DENTRO del documento principal.
% No contiene documentclass, usepackage, begin{document}, end{document} ni appendix.
% Usa TikZ, booktabs y los colores upmblue/upmcyan/upmorange/upmgray/upmlight,
% que ya estan definidos en vuestra plantilla.
%
% En vuestro main.tex: reemplazar TODO el bloque desde
%   \section[Enfoques híbridos]{Enfoques híbridos y líneas futuras}
% hasta justo ANTES de \begin{frame}[t]{Referencias}
% por el contenido completo de este archivo.
%
% Duracion prevista: 14:00 de contenido + ~0:30 de apertura y transicion,
% con margen para ajustarse a los 15 minutos de la seccion.
% ============================================================================

% --- Estilos específicos de la sección de enfoques híbridos ---
\tikzset{
  hybridcard/.style={draw=upmblue!45, fill=upmlight, rounded corners=4pt,
    align=center, line width=.7pt, minimum height=1.50cm,
    inner sep=5pt, font=\small},
  hybridstep/.style={draw=upmblue!75, fill=upmblue!7, rounded corners=3pt,
    align=center, minimum height=1.25cm, text width=2.45cm,
    inner sep=5pt, font=\small},
  hybridorange/.style={draw=upmorange!80, fill=upmorange!7, rounded corners=3pt,
    align=center, minimum height=1.25cm, text width=2.45cm,
    inner sep=5pt, font=\small},
  hybridarrow/.style={-{Latex[length=2.3mm]}, draw=upmorange, line width=1.2pt}
}
\providecommand{\hybridsource}[1]{\par\vfill{\color{upmgray}\fontsize{7.6}{9}\selectfont\hfill #1}}
\providecommand{\hybridtag}[1]{\textcolor{upmcyan}{\bfseries #1}}

% Macros añadidas solo para los gráficos del bloque híbrido.
\tikzset{
  hybridvertex/.style={circle,draw=upmblue,fill=upmlight,
    minimum size=.69cm,inner sep=1pt,font=\small\bfseries},
  hybridedge/.style={draw=upmblue,line width=1.4pt},
  hybridstat/.style={draw=upmcyan!70,fill=upmcyan!7,rounded corners=4pt,
    minimum height=1.36cm,align=center,inner sep=6pt,font=\normalsize}
}

\section[Enfoques híbridos]{Enfoques híbridos y líneas futuras}

% 01 / 00:55 — Motivación.
% GUION PARA QUIEN PRESENTE: Explicar que constraint-based y score-based resuelven problemas distintos y sus debilidades motivan combinarlos.
\begin{frame}[t]{Motivación}
  \vspace{.12cm}
  \begin{columns}[T,onlytextwidth]
    \begin{column}{.48\textwidth}
      \begin{block}{\textit{Constraint-based}}
        \begin{itemize}\setlength{\itemsep}{5pt}
          \item Pruebas de \textbf{independencia condicional} (IC).
          \item El número de pruebas puede crecer mucho.
          \item Un error temprano puede \textbf{propagarse}.
        \end{itemize}
      \end{block}
    \end{column}
    \begin{column}{.48\textwidth}
      \begin{block}{\textit{Score-and-search}}
        \begin{itemize}\setlength{\itemsep}{5pt}
          \item Asigna una \textbf{puntuación} a cada estructura.
          \item El espacio de DAGs es enorme.
          \item La búsqueda voraz puede acabar en \textbf{óptimos locales}.
        \end{itemize}
      \end{block}
    \end{column}
  \end{columns}
  \vspace{.20cm}
  \begin{exampleblock}{La idea híbrida}
    \centering \textbf{Usar las dependencias entre variables para orientar la optimización.}\\
    Combinar fortalezas \emph{no} garantiza superar a ambos métodos por separado.
  \end{exampleblock}
  \hybridsource{Fuente: Kitson et al. (2023)}
\end{frame}

% 02 / 01:10 — Arquitectura.
% GUION PARA QUIEN PRESENTE: Describir restrict como filtro y maximise como optimizacion; recordar el riesgo de eliminar una arista cierta.
\begin{frame}[t]{\textit{Restrict/Maximise}: ¿cómo funciona un híbrido clásico?}
  \vspace{.13cm}
  \begin{center}
    \begin{tikzpicture}[node distance=.37cm]
      \node[hybridstep,text width=2.3cm,minimum height=1.55cm] (datos) {\textbf{Datos}\\[3pt]\small Observaciones de\\las variables};
      \node[hybridstep,text width=2.50cm,minimum height=1.55cm,right=of datos] (filtro) {\textbf{1. RESTRICT}\\[3pt]\small Tests de IC:\\candidatos};
      \node[hybridorange,text width=2.50cm,minimum height=1.55cm,right=of filtro] (opt) {\textbf{2. MAXIMISE}\\[3pt]\small Búsqueda por\\puntuación};
      \node[hybridstep,text width=2.2cm,minimum height=1.55cm,right=of opt] (out) {\textbf{Resultado}\\[3pt]\small DAG con aristas\\orientadas};
      \draw[hybridarrow] (datos)--(filtro);
      \draw[hybridarrow] (filtro)--(opt);
      \draw[hybridarrow] (opt)--(out);
    \end{tikzpicture}
  \end{center}
  \vspace{.24cm}
  \begin{columns}[T,onlytextwidth]
    \begin{column}{.48\textwidth}
      \begin{block}{Ventaja esperada}
        \centering Menos candidatos \(\Rightarrow\) búsqueda más manejable.
      \end{block}
    \end{column}
    \begin{column}{.48\textwidth}
      \begin{alertblock}{Riesgo central}
        \centering La búsqueda no recupera una arista \textbf{prohibida por el filtro}.
      \end{alertblock}
    \end{column}
  \end{columns}
  \hybridsource{Fuente: Kitson et al. (2023). Esquema conceptual, no algoritmo específico.}
\end{frame}

% 03 / 01:10 — Ejemplo pedagogico; figura deliberadamente ficticia.
% GUION PARA QUIEN PRESENTE: El dibujo es inventado: si la arista B-D queda fuera, ningun score la puede recuperar.
\begin{frame}[t]{Ejemplo: una restricción demasiado estricta}
  \begin{columns}[T,onlytextwidth]
    \begin{column}{.47\textwidth}
      \begin{block}{Esqueleto que querríamos recuperar}
        \centering
        \begin{tikzpicture}[scale=.91,every node/.style={transform shape}]
          \node[hybridvertex] (a) at (0,1.5) {A};
          \node[hybridvertex] (b) at (2.5,1.5) {B};
          \node[hybridvertex] (c) at (0,0) {C};
          \node[hybridvertex] (d) at (2.5,0) {D};
          \draw[hybridedge] (a)--(b);
          \draw[hybridedge] (b)--(d);
          \draw[hybridedge] (c)--(d);
        \end{tikzpicture}
        \vspace{.09cm}
        \par\small Tres conexiones relevantes.
      \end{block}
    \end{column}
    \begin{column}{.47\textwidth}
      \begin{block}{Candidatos tras un test erróneo}
        \centering
        \begin{tikzpicture}[scale=.91,every node/.style={transform shape}]
          \node[hybridvertex] (a) at (0,1.5) {A};
          \node[hybridvertex] (b) at (2.5,1.5) {B};
          \node[hybridvertex] (c) at (0,0) {C};
          \node[hybridvertex] (d) at (2.5,0) {D};
          \draw[hybridedge] (a)--(b);
          \draw[hybridedge] (c)--(d);
          \draw[upmorange,densely dashed,line width=1.5pt] (b)--(d);
          \node[font=\scriptsize,text=upmorange,anchor=west] at (2.65,.65) {excluida};
        \end{tikzpicture}
        \vspace{.09cm}
        \par\small La arista B--D queda fuera de la búsqueda.
      \end{block}
    \end{column}
  \end{columns}
  \vspace{.20cm}
  \begin{alertblock}{Problema de fondo}
    \centering Optimizar muy bien \emph{dentro} del espacio permitido no corrige un filtro equivocado.
  \end{alertblock}
  \hybridsource{Ejemplo didáctico inventado para explicar el compromiso; no corresponde a un experimento.}
\end{frame}

% 04 / 01:05 — CB.
% GUION PARA QUIEN PRESENTE: CB usa informacion de PC para obtener un orden; K2 usa ese orden y las etapas se alternan.
\begin{frame}[t]{CB (1993): restringir mediante una ordenación}
  \vspace{.14cm}
  \begin{center}
    \begin{tikzpicture}[node distance=.56cm]
      \node[hybridstep,text width=2.9cm,minimum height=1.25cm] (pc) {\textbf{PC + tests de IC}\\\small Información estructural};
      \node[hybridstep,text width=2.9cm,minimum height=1.25cm,right=of pc] (ord) {\textbf{Orden de variables}\\\small Limita padres posibles};
      \node[hybridorange,text width=2.9cm,minimum height=1.25cm,right=of ord] (k2) {\textbf{K2}\\\small Búsqueda voraz con score};
      \draw[hybridarrow] (pc)--(ord);
      \draw[hybridarrow] (ord)--(k2);
      \draw[hybridarrow,rounded corners=6pt] (k2.south) -- ++(0,-.85) -| node[pos=.25,below,font=\scriptsize,text=upmorange] {Itera y amplía conjuntos condicionantes} (pc.south);
    \end{tikzpicture}
  \end{center}
  \vspace{.35cm}
  \begin{columns}[T,onlytextwidth]
    \begin{column}{.48\textwidth}
      \begin{block}{Su aportación}
        \small Combinar información de PC con K2, alternando ambas etapas.
      \end{block}
    \end{column}
    \begin{column}{.48\textwidth}
      \begin{alertblock}{Su límite}
        \small K2 depende del orden obtenido y no garantiza un óptimo global.
      \end{alertblock}
    \end{column}
  \end{columns}
  \hybridsource{Fuente: Singh y Valtorta (1993)}
\end{frame}

% 05 / 01:05 — Sparse Candidate.
% GUION PARA QUIEN PRESENTE: Sparse Candidate cambia un orden global por listas locales de padres y actualiza candidatos iterativamente.
\begin{frame}[t]{Sparse Candidate (1999): restringir posibles padres}
  \vspace{.05cm}
  \begin{columns}[T,onlytextwidth]
    \begin{column}{.43\textwidth}
      \begin{block}{Para cada variable objetivo \(X\)}
        \vspace{.09cm}
        \centering
        \begin{tikzpicture}
          \node[hybridvertex,minimum size=.66cm] (x) at (1.45,1.00) {X};
          \node[hybridvertex,minimum size=.56cm] (a) at (0,2.0) {A};
          \node[hybridvertex,minimum size=.56cm] (b) at (2.9,2.0) {B};
          \node[hybridvertex,minimum size=.56cm] (c) at (0,0) {C};
          \node[hybridvertex,minimum size=.56cm] (d) at (2.9,0) {D};
          \draw[hybridarrow] (a)--(x);
          \draw[hybridarrow] (b)--(x);
          \draw[upmgray,dashed,thin] (c)--(x);
          \draw[upmgray,dashed,thin] (d)--(x);
        \end{tikzpicture}
        \par\vspace{.06cm}
        \scriptsize Azul = candidatos; gris = excluidos. Ejemplo ilustrativo.
      \end{block}
    \end{column}
    \begin{column}{.53\textwidth}
      \begin{enumerate}\setlength{\itemsep}{9pt}
        \item \textbf{Seleccionar} un pequeño conjunto de padres candidatos mediante dependencias estadísticas.
        \item \textbf{Optimizar} el DAG mediante Hill-Climbing con lista tabú sobre esos candidatos.
        \item \textbf{Repetir}: actualizar candidatos con la estructura aprendida.
      \end{enumerate}
    \end{column}
  \end{columns}
  \vspace{.13cm}
  \begin{exampleblock}{Salto frente a CB}
    \centering La restricción deja de ser principalmente un \textbf{orden global}: pasa a ser una \textbf{lista local de padres posibles}.
  \end{exampleblock}
  \hybridsource{Fuente: Friedman et al. (1999)}
\end{frame}

% 06 / 01:15 — MMHC.
% GUION PARA QUIEN PRESENTE: MMHC usa MMPC para detectar padres e hijos, luego HC para aprender el DAG sin volver a alternar.
\begin{frame}[t]{MMHC (2006): separar detección y optimización}
  \vspace{.09cm}
  \begin{center}
    \begin{tikzpicture}[node distance=.6cm]
      \node[hybridstep,text width=3.6cm,minimum height=1.52cm] (mmpc) {\textbf{1. MMPC}\\[3pt]\small Tests de IC para detectar\\posibles padres e hijos};
      \node[hybridorange,text width=3.6cm,minimum height=1.52cm,right=of mmpc] (hc) {\textbf{2. Hill-Climbing}\\[3pt]\small Puntúa, selecciona y orienta\\aristas permitidas};
      \draw[hybridarrow] (mmpc)--(hc);
      \node[below=.10cm of mmpc,font=\scriptsize,text=upmgray] {Estimación del esqueleto};
      \node[below=.10cm of hc,font=\scriptsize,text=upmgray] {Construcción de un DAG};
    \end{tikzpicture}
  \end{center}
  \vspace{.58cm}
  \begin{columns}[T,onlytextwidth]
    \begin{column}{.48\textwidth}
      \begin{block}{Respecto a Sparse Candidate}
        \small Candidatos locales adaptados a cada variable; \textbf{dos etapas}, sin alternancia iterativa.
      \end{block}
    \end{column}
    \begin{column}{.48\textwidth}
      \begin{alertblock}{Dependencia que persiste}
        \small Si MMPC omite una relación, Hill-Climbing no puede añadirla.
      \end{alertblock}
    \end{column}
  \end{columns}
  \hybridsource{Fuente: Tsamardinos et al. (2006)}
\end{frame}

% 07 / 01:00 — Comparación de tres enfoques.
% GUION PARA QUIEN PRESENTE: La tabla es un recapitulativo, no es necesario leerla entera: enfocar que cambia el tipo de restriccion.
\begin{frame}[t]{Tres hitos: lo que realmente cambia}
  \vspace{.10cm}
  \small
  \renewcommand{\arraystretch}{1.55}
  \setlength{\tabcolsep}{4pt}
  \begin{tabular}{@{}p{.18\textwidth}p{.33\textwidth}p{.20\textwidth}p{.22\textwidth}@{}}
    \toprule
    \textbf{Método} & \textbf{¿Qué restringe?} & \textbf{Búsqueda} & \textbf{Organización} \\
    \midrule
    \textbf{CB} (1993) & Ordenación de variables & K2 & Iterativa \\
    \textbf{Sparse Candidate} (1999) & Posibles padres por nodo & HC + tabú & Iterativa \\
    \textbf{MMHC} (2006) & Adyacencias padre--hijo (MMPC) & HC & Dos etapas \\
    \bottomrule
  \end{tabular}
  \vspace{.42cm}
  \begin{exampleblock}{Hilo histórico}
    \centering \textbf{De guiar la búsqueda con un orden a delimitar las relaciones candidatas.}\\[4pt]
    \small A medida que mejora el filtro, resulta más importante saber \emph{qué ocurre cuando falla}.
  \end{exampleblock}
  \hybridsource{Fuentes: Singh y Valtorta (1993); Friedman et al. (1999); Tsamardinos et al. (2006)}
\end{frame}

% 08 / 01:05 — H2PC.
% GUION PARA QUIEN PRESENTE: H2PC intenta evitar falsos negativos en el filtro MMPC, a cambio de coste computacional.
\begin{frame}[t]{H2PC (2014): evitar falsos negativos en el filtro}
  \vspace{.08cm}
  \begin{columns}[T,onlytextwidth]
    \begin{column}{.45\textwidth}
      \begin{block}{La modificación}
        \vspace{.06cm}
        \centering
        \begin{tikzpicture}[node distance=.38cm]
          \node[hybridstep,text width=3.7cm,minimum height=.95cm] (mm) {\textbf{MMHC: MMPC}\\\small Candidatos padre--hijo};
          \node[hybridorange,text width=3.7cm,minimum height=.95cm,below=of mm] (hp) {\textbf{H2PC: HPC}\\\small Recuperación local más sensible};
          \draw[hybridarrow] (mm)--(hp);
        \end{tikzpicture}
      \end{block}
    \end{column}
    \begin{column}{.51\textwidth}
      \begin{itemize}\setlength{\itemsep}{9pt}
        \item Falso negativo: \textbf{una arista real no pasa el filtro}.
        \item HPC combina técnicas locales para mejorar la detección.
        \item La etapa \textit{score-based} sigue buscando con los candidatos.
      \end{itemize}
    \end{column}
  \end{columns}
  \vspace{.21cm}
  \begin{alertblock}{El nuevo compromiso}
    \centering \textbf{Menos aristas reales descartadas} puede implicar \textbf{más tiempo de ejecución}.
  \end{alertblock}
  \hybridsource{Fuente: Gasse et al. (2014). Mejora observada frente a MMHC en sus experimentos.}
\end{frame}

% 09 / 01:05 — ARGES.
% GUION PARA QUIEN PRESENTE: ARGES flexibiliza las restricciones durante la busqueda por clases de equivalencia, distinto de mejorar el filtro.
\begin{frame}[t]{ARGES (2018): que las restricciones sean adaptativas}
  \vspace{.10cm}
  \begin{columns}[T,onlytextwidth]
    \begin{column}{.48\textwidth}
      \begin{block}{Una búsqueda restringida rígida}
        \centering\vspace{.15cm}
        \textcolor{upmblue}{\large\bfseries Esqueleto inicial}\\[10pt]
        \textcolor{upmorange}{\Large$\Downarrow$}\\[5pt]
        \small Solo se permiten movimientos compatibles con ese filtro.
        \vspace{.15cm}
      \end{block}
    \end{column}
    \begin{column}{.48\textwidth}
      \begin{block}{ARGES: restricciones adaptativas}
        \centering\vspace{.15cm}
        \textcolor{upmblue}{\large\bfseries Filtro inicial}\\[10pt]
        \textcolor{upmorange}{\Large$\Downarrow$}\\[5pt]
        \small Admite \textbf{algunas aristas adicionales} según la configuración actual del grafo.
        \vspace{.15cm}
      \end{block}
    \end{column}
  \end{columns}
  \vspace{.20cm}
  \begin{exampleblock}{Diferencia respecto a H2PC}
    \centering H2PC mejora \textbf{qué relaciones detectamos}.\\
    ARGES modifica \textbf{cómo aplicamos las restricciones} durante la búsqueda.
  \end{exampleblock}
  \hybridsource{Fuente: Nandy et al. (2018). Consistencia demostrada bajo supuestos específicos.}
\end{frame}

% 10 / 01:05 — Alternativas.
% GUION PARA QUIEN PRESENTE: GFCI invierte el orden y trata posibles latentes; SaiyanH impone conectividad y ese supuesto puede fallar.
\begin{frame}[t]{Otras formas de combinar ambos paradigmas}
  \vspace{.08cm}
  \begin{columns}[T,onlytextwidth]
    \begin{column}{.48\textwidth}
      \begin{block}{GFCI (2016): invertir el orden}
        \centering \textbf{Score (FGS)} $\longrightarrow$ \textbf{IC / FCI}
        \vspace{.17cm}
        \begin{itemize}\setlength{\itemsep}{5pt}
          \item Empieza con una estructura aprendida por puntuación.
          \item La refina mediante tests y orientación causal.
          \item Devuelve un \textbf{PAG}, admite posibles variables latentes.
        \end{itemize}
      \end{block}
    \end{column}
    \begin{column}{.48\textwidth}
      \begin{block}{SaiyanH (2020): añadir un supuesto}
        \centering \textbf{Grafo conectado} $\longrightarrow$ \textbf{IC + score}
        \vspace{.17cm}
        \begin{itemize}\setlength{\itemsep}{5pt}
          \item Combina asociaciones, IC y búsqueda HC con tabú.
          \item Conserva la \textbf{conectividad} del grafo.
          \item Si el supuesto es falso, puede crear aristas espurias.
        \end{itemize}
      \end{block}
    \end{column}
  \end{columns}
  \vspace{.12cm}
  \centering\textcolor{upmblue}{\bfseries ``Híbrido'' no significa siempre ``filtrar y después buscar''.}
  \hybridsource{Fuentes: Ogarrio et al. (2016); Constantinou (2020)}
\end{frame}

% 11 / 01:00 — Recientes.
% GUION PARA QUIEN PRESENTE: ESLH intenta abaratar tests y SAPSO intenta explorar mejor que el HC voraz; los resultados son contextuales.
\begin{frame}[t]{Avances recientes: acelerar el filtro o mejorar la búsqueda}
  \vspace{.15cm}
  \begin{columns}[T,onlytextwidth]
    \begin{column}{.48\textwidth}
      \begin{block}{ESLH (2024)}
        \centering\hybridtag{Coste de las pruebas de IC}
        \vspace{.13cm}
        \begin{itemize}\setlength{\itemsep}{6pt}
          \item Umbrales dinámicos y eliminación de triángulos.
          \item Aprendizaje más eficiente del esqueleto.
          \item \textbf{HC en tres etapas} menos sujeto a restricciones rígidas.
        \end{itemize}
      \end{block}
    \end{column}
    \begin{column}{.48\textwidth}
      \begin{block}{SAPSO (2023)}
        \centering\hybridtag{Exploración del espacio de DAGs}
        \vspace{.13cm}
        \begin{itemize}\setlength{\itemsep}{6pt}
          \item Información mutua para guiar y restringir.
          \item Búsqueda poblacional \textbf{PSO} con score BIC.
          \item Reinicio adaptativo cuando la población se estanca.
        \end{itemize}
      \end{block}
    \end{column}
  \end{columns}
  \vspace{.00cm}
  \begin{alertblock}{Interpretar bien los resultados}
    \centering Las mejoras son experimentales y contextuales; \textbf{SAPSO se evaluó en dos redes}.
  \end{alertblock}
  \hybridsource{Fuentes: Wang et al. (2024); Liu et al. (2023)}
\end{frame}

% 12 / 01:10 — Evidencia comparativa.
% GUION PARA QUIEN PRESENTE: No hay superioridad universal: citar el alcance limitado del estudio comparativo.
\begin{frame}[t]{¿Son siempre mejores los híbridos? La evidencia dice que no}
  \vspace{.17cm}
  \begin{center}
    \begin{tikzpicture}[node distance=.45cm]
      \node[hybridstat,text width=2.52cm] (a) {\textbf{14}\\[3pt]\small algoritmos\\comparados};
      \node[hybridstat,text width=2.52cm,right=of a] (b) {\textbf{5}\\[3pt]\small estructuras\\simuladas};
      \node[hybridstat,text width=3.12cm,right=of b] (c) {\textbf{MCC, F1, SHD}\\[3pt]\small varias métricas y\\tamaños muestrales};
    \end{tikzpicture}
  \end{center}
  \vspace{.35cm}
  \begin{alertblock}{Resultado del estudio comparativo de Karadag (2026)}
    \centering \textbf{Ninguna familia dominó sistemáticamente} en las condiciones estudiadas.\\[4pt]
    \small MMHC y H2PC no superaron en todos los casos a enfoques puramente \textit{constraint-based}.
  \end{alertblock}
  \vspace{.16cm}
  \centering\small La comparación debe considerar \textbf{topología, muestras, supuestos y métrica}.
  \hybridsource{Fuente: Karadag (2026). La conclusión solo se refiere a los escenarios evaluados.}
\end{frame}

% 13 / 00:55 — Cierre.
% GUION PARA QUIEN PRESENTE: Cerrar resumiendo las tres preguntas clave: que restringir, como buscar y bajo que supuestos.
\begin{frame}[t]{Conclusiones y líneas futuras del bloque}
  \vspace{.14cm}
  \begin{columns}[T,onlytextwidth]
    \begin{column}{.315\textwidth}
      \begin{block}{1. ¿Qué restringir?}
        \small Evitar falsos negativos con pocos datos; filtrar sin perder relaciones relevantes.\\[4pt]
        \textcolor{upmcyan}{\bfseries H2PC, ESLH}
      \end{block}
    \end{column}
    \begin{column}{.315\textwidth}
      \begin{block}{2. ¿Cómo buscar?}
        \small Explorar mejor DAGs o clases de equivalencia; aplicar restricciones con flexibilidad.\\[4pt]
        \textcolor{upmcyan}{\bfseries ARGES, SAPSO}
      \end{block}
    \end{column}
    \begin{column}{.315\textwidth}
      \begin{block}{3. ¿Bajo qué supuestos?}
        \small Tratar variables latentes y evaluar métodos en escenarios variados.\\[4pt]
        \textcolor{upmcyan}{\bfseries GFCI, evaluación}
      \end{block}
    \end{column}
  \end{columns}
  \vspace{.27cm}
  \begin{exampleblock}{Idea para cerrar la sección}
    \centering \large\textbf{El reto no es simplemente mezclar técnicas: es decidir cómo integrarlas y cuándo convienen.}
  \end{exampleblock}
  \hybridsource{Síntesis de la revisión del estado del arte (Kitson et al., 2023, y trabajos citados).}
\end{frame}


% ---------------------------------------------------------------------------
\begin{frame}[t]{Conclusiones}
  \begin{itemize}
    \item \slot{Conclusión 1}
    \item \slot{Conclusión 2}
    \item \slot{Conclusión 3}
    \item \slot{Conclusión 4}
  \end{itemize}
  \vfill
  {\scriptsize\color{upmgray}% TODO: sustituir por las conclusiones propias del informe.
  [Esqueleto: completar con las conclusiones del informe.]}
\end{frame}

% ---------------------------------------------------------------------------
\begin{frame}[t]{Líneas futuras}
  \begin{columns}[T,onlytextwidth]
    \begin{column}{0.48\textwidth}
      \begin{block}{Línea 1}
        \slot{Descripción breve}
      \end{block}
      \begin{block}{Línea 2}
        \slot{Descripción breve}
      \end{block}
    \end{column}
    \begin{column}{0.49\textwidth}
      \begin{block}{Línea 3}
        \slot{Descripción breve}
      \end{block}
      \begin{block}{Línea 4}
        \slot{Descripción breve}
      \end{block}
    \end{column}
  \end{columns}
\end{frame}

% ---------------------------------------------------------------------------
\begin{frame}[t]{Referencias: algoritmo PC y variantes}
  \scriptsize
  \begin{itemize}\setlength{\itemsep}{1pt}
    \item Spirtes, P., Glymour, C. (1991). An algorithm for fast recovery of sparse causal graphs. \emph{Social Science Computer Review}.
    \item Meek, C. (1995). Causal inference and causal explanation with background knowledge. \emph{UAI}.
    \item Ramsey, J., Zhang, J., Spirtes, P. (2006). Adjacency-faithfulness and conservative causal inference. \emph{UAI}.
    \item Abellán, J., Gómez-Olmedo, M., Moral, S. (2006). Some variations on the PC algorithm. \emph{PGM}.
    \item Kalisch, M., Bühlmann, P. (2007). Estimating high-dimensional DAGs with the PC-algorithm. \emph{JMLR}.
    \item Uhler, C., Raskutti, G., Bühlmann, P., Yu, B. (2013). Geometry of the faithfulness assumption. \emph{Annals of Statistics}.
    \item Colombo, D., Maathuis, M. (2014). Order-independent constraint-based causal structure learning. \emph{JMLR}.
    \item Madsen, A.\,L. et al.\ (2017). A parallel algorithm for BN structure learning from large data sets. \emph{KBS}.
    \item Scutari, M. (2017). BN constraint-based structure learning: parallel and optimized implementations in bnlearn. \emph{JSS}.
    \item Le, T.\,D. et al.\ (2019). A fast PC algorithm for high dimensional causal discovery with multi-core PCs. \emph{IEEE/ACM TCBB}.
    \item Runge, J. et al.\ (2019). Detecting and quantifying causal associations in large nonlinear time series datasets. \emph{Science Advances}.
    \item Kocaoglu, M. (2023). Characterization and learning of causal graphs with small conditioning sets. \emph{NeurIPS}.
    \item Biswas, R., Mukherjee, S. (2024). Consistent causal inference from time series with PC algorithm. \emph{Statistics and Computing}.
    \item Zhou, J. et al.\ (2025). Differentiable constraint-based causal discovery. \emph{NeurIPS}.
    \item Améndola, C., Hollering, B., Nowell, F. (2026). A PC algorithm for max-linear Bayesian networks. \emph{JCGS}.
  \end{itemize}
  \vfill
  {\scriptsize\color{upmgray}La bibliografía completa está en el informe.}
\end{frame}

% ---------------------------------------------------------------------------
\begin{frame}[plain]
  \begin{tikzpicture}[remember picture,overlay]
    \fill[upmblue] (current page.south west) rectangle (current page.north east);
    \node[anchor=center,align=center] at (current page.center) {%
      {\color{white}\fontsize{34}{38}\selectfont\bfseries Gracias}\\[0.35cm]
      {\color{white!80}\large Preguntas y discusión}\\[0.6cm]
      {\color{upmorange}\rule{4cm}{3pt}}\\[0.5cm]
      {\color{white!70}\small Damián Benasco \quad Thales Camargo \quad Juan David Cárdenas}
    };
  \end{tikzpicture}
\end{frame}

% ===========================================================================
%  ANEXO
% ===========================================================================
\appendix

\begin{frame}[t]{Anexo: material de respaldo}
  \begin{itemize}
    \item \slot{Tabla de variantes del PC}
    \item \slot{Figura de decisión sobre contrastes}
    \item \slot{Detalles que no caben en el cuerpo principal}
  \end{itemize}
\end{frame}

% ---------------------------------------------------------------------------
\begin{frame}[t]{Anexo: cuándo elegir cada contraste}
  \footnotesize
  \setlength{\tabcolsep}{3pt}
  \renewcommand{\arraystretch}{1.10}
  \begin{tabular}{@{}>{\raggedright\arraybackslash}p{0.21\textwidth}
                     >{\raggedright\arraybackslash}p{0.45\textwidth}
                     >{\raggedright\arraybackslash}p{0.27\textwidth}@{}}
    \toprule
    \textbf{Modelo} & \textbf{Cuándo usarlo} & \textbf{Coste y límite} \\
    \midrule
    Lineal-gaussiano
      & continuas, normalidad conjunta, linealidad; $n$ grande y
        $|\mathbf Z|$ moderado
      & $\mathcal O(|\mathbf Z|^3)$, $p$ analítico; ciego a $Y=X^2$ \\
    \addlinespace
    Medias condicionales
      & continuas o mixtas; media condicional especificable (GLM, doble ML)
      & 2 regresiones $+\,\mathcal O(n)$; WGCM si la varianza del error no es constante \\
    \addlinespace
    Sin forma funcional
      & no lineal o error con varianza no constante; $n$ y $|\mathbf Z|$ moderados
      & KCI $\mathcal O(n^3)$; RCIT/RCoT $\mathcal O(n)$; frágil si $n<1000$ \\
    \addlinespace
    Tabla de contingencia
      & categóricas; conteos por celda suficientes; $|\mathbf Z|$ bajo
      & $\chi^2/G^2$ analítico; SECMI si $|\mathbf Z|$ grande \\
    \bottomrule
  \end{tabular}
  \vspace{0.1cm}
  \begin{alertblock}{GCM frente a KCI}
    GCM si la dependencia pasa por la media condicional y el regresor es
    fiable; KCI si no quieres imponer esa forma y asumes su coste.
  \end{alertblock}
\end{frame}

\end{document}
